\section{Introduction}\label{sec:intro}

Optimization models in process systems engineering, energy networks and
optimal control are often built from local pieces. A flowsheet links unit
models along streams, a multiperiod model links consecutive time stages, and a
network model links quantities at adjacent nodes. The objective is then a sum
of terms that each depend on a few variables, and the graph recording which
variables share a term often has small treewidth. For variables with finitely
many values, small treewidth is exactly what dynamic programming needs:
nonserial dynamic programming, bucket elimination and junction-tree message
passing find a global optimum in time exponential only in the width of a tree
decomposition \cite{BerteleBrioschi1972,Dechter1999,KollerFriedman2009}.

Mixed-integer nonlinear models have continuous variables and nonconvex
terms, and there the situation is less clear. Global solvers based on spatial
branch-and-bound bound individual terms by convex relaxations and branch on
variable domains \cite{TawarmalaniSahinidis2005,BelottiEtAl2009,
MisenerFloudas2014,BelottiEtAl2013}; they do not use a tree decomposition to
organize the search. Discretizing the continuous variables and running a
dynamic program gives neither a valid bound nor a running time that scales
well with the accuracy: a uniform grid of mesh $h$ has $\Theta(1/h)$ points per
coordinate, so tables grow like $h^{-p}$ in the bag size $p$. Sparse polynomial
optimization exploits the same structure through linear and semidefinite
hierarchies \cite{WakiEtAl2006,Lasserre2006,BienstockMunoz2018}, again with
sizes polynomial in the inverse accuracy.

Small treewidth alone does not make continuous nonconvex problems
tractable. Del Pia and Khajavirad \cite{DelPiaKhajavirad2026} recently showed
that box-constrained nonconvex quadratic programming is solvable in strongly
polynomial time on forests, but strongly NP-hard at treewidth two even with
small integer coefficients, and that quartic minimization is strongly NP-hard
on paths. This paper asks which additional, checkable structure makes tree
decompositions useful for \emph{certified} global optimization of
mixed-integer problems, and what fails when it is absent.

\subsection{Main results}
We consider $\min\{F(x):x\in X\}$, where $F=\sum_af_a(x_{S_a})$ is a sum of
factors, $X$ is a box in which some coordinates are integer, and a tree
decomposition of the scopes $S_a$ with maximum bag size $p$ is given. Two
quantities govern the results: upper bounds $L_i$ on the curvature of $F$ along
each coordinate (the diagonal of the Hessian for a quadratic), and quadratic
growth $F(x)-\OPT\ge g\norm{x-x^*}^2$ at a unique minimizer $x^*$. We write
$\kappa=\max\{1,L/g\}$ with $L=\max_iL_i$, and $\bar\kappa\le\kappa$ for a
weighted version that does not change when coordinates are rescaled
(Definition~\ref{def:growth}).

\paragraph{Certified coordinate grids.}
Choose finite grids in the coordinates and subtract, at each node $v$ of
coordinate $i$, the unary correction $L_iw^2/8$, where $w$ is the largest grid
interval ending at $v$. The minimum of the corrected objective over the grid is
a lower bound on $\min_XF$; in fact it equals the best of the vertex bounds of
Bajaj and Hasan \cite{BajajHasan2020} over all cells of the product grid
(Proposition~\ref{prop:cellwise}). Because the corrections are unary, this
minimum and all coordinate min-marginals are computed by ordinary dynamic
programming on the given decomposition, without enumerating cells. The
min-marginals discard coordinate intervals that cannot contain a point better
than a known feasible value, and the grids of the successive stages, one
bound and one feasible point form a certificate that an independent program
checks by recomputation (Theorem~\ref{thm:certificate}). None of this uses
convexity, uniqueness or growth.

\paragraph{Accuracy-independent grids and exact output.}
Under weighted growth, grids whose spacing grows linearly away from the
current corrected minimizer, combined with this filtering, keep
$O(\sqrt{\bar\kappa}\log n)$ nodes per coordinate at every stage, independently of
the accuracy (Lemma~\ref{lem:states}); capped trials remove the need to know the
growth constant. For an explicit rational mixed-integer quadratic program on a
box this gives a certified $2^{-q}$-approximation in
$f(p,\bar\kappa)(I+q+1)^5$ bit operations, where $I$ is the input length
(Theorem~\ref{thm:approx}), and the exact minimizer and optimal value in
$f(p,\kappa)I^{O(1)}$ bit operations (Theorem~\ref{thm:exact}). The algorithms
are therefore fixed-parameter tractable in the pair of bag size and condition
number. A snapping procedure based on explicit rational height bounds turns
any point within $\poly(I)+\log_2(1/g_S)$ bits of the optimum into an exact
minimizer, where $g_S$ is the growth constant towards the optimal \emph{set},
so exactness itself needs no uniqueness (Theorem~\ref{thm:transfer}). The
approximation result extends to fixed-degree polynomial factors, with exact
output when all coordinates are integer (Corollary~\ref{cor:poly}).

A chain with box widths $2^t$ (Example~\ref{ex:chain}) shows the difference from
exact parametric dynamic programming: its exact Bellman messages need at least
$2^{m-1}$ quadratic pieces, while the corrected grids certify any accuracy in
time polynomial in $m$. Our implementation certifies a gap of $10^{-3}$ for
$m=64$ with at most eleven nodes per coordinate grid
(Section~\ref{sec:computation}).

\paragraph{Conditional recourse.}
The condition number uses diagonal curvature, so a stiff convex coupling term
inflates it although convex subproblems are easy. The corrected-grid bound
applies to any value function $V(v)=\min_yF(v,y)$ whose inner feasible set does
not depend on $v$, because partial minimization preserves coordinate curvature
and growth (Lemma~\ref{lem:valuefunction}); bag-local corrections are valid
only with such exact or certified recourse (Theorem~\ref{thm:cr-filter}), not
with grid messages (Proposition~\ref{prop:star}). Private convex quadratic
blocks become value factors whose curvature is certified by piecewise-affine
responses; the certified cancellation is the best possible for the block, and
no global overlay of the pieces is needed (Theorem~\ref{thm:cv}). A large
residual problem that is concave in each coordinate and submodular after sign
changes is removed from discretization, because each conditional value is one
minimum cut; only a small core is searched (Section~\ref{sec:cuts}). For balanced
quadratics the grid problem itself is a minimum cut, which gives certified
approximation in time polynomial in $n$, $\kappa$ and $\log(1/\varepsilon)$
without any width parameter (Theorem~\ref{thm:balanced}). Recourse does not
reduce the parameter to negative curvature over growth in general
(Proposition~\ref{prop:cv-limit}).

\paragraph{Constraints and several minimizers.}
Independent rounding, on which the certificate rests, breaks under coupling
constraints. When the continuous coupling matrix is totally unimodular and the
data are aligned to a mesh, correlated rounding to feasible cell corners
restores the certificate, with exact feasibility and exact output for
quadratics; the algorithm is polynomial for fixed width and conditioning but
not fixed-parameter tractable, and we show that the uniform meshes responsible
for this are forced (Section~\ref{sec:constraints}). With several minimizers, a
single graded grid can need exponential time even for two isolated minimizers
(Proposition~\ref{prop:twocenters}); unions of uniform cells restore a
polynomial bound for fixed width when the optimal set is finite
(Theorem~\ref{thm:cells}). For coordinatewise concave mixed-integer
quadratics, the Bellman residuals of an endpoint dynamic program describe the
entire optimal set as a tree-structured constraint problem
(Corollary~\ref{cor:facecsp}), and for instances certified by a diagonal
Lagrangian multiplier an unknown optimal set is found exactly in
$f(p,\kappa)\poly(I)$ time (Theorem~\ref{thm:diagdiscovery}).

\paragraph{Structural limits.}
Section~\ref{sec:limits} shows that each hypothesis does real work. Box QP with
bag size three and a unique minimizer is NP-hard on instances with
$\log\kappa=O(I)$, so the dependence on $\kappa$ cannot be polylogarithmic unless
P${}={}$NP; under the exponential-time hypothesis the dependence on $p$ is
exponential even for $\kappa\le2$; under its randomized version the exponent of
$\kappa$ must grow linearly with $p$; and in the value-oracle model the
exponent $p/2$ of $\kappa$ is optimal. Further results show that growth towards
a continuum of minimizers does not bound any corrected-grid certificate, that
one local affine equality chain makes the problem NP-hard at $\kappa=1$, and
that matching separator moments of any finite order does not give a local
error bound in terms of negative curvature.

\subsection{Contributions}
The ingredients we use are classical: interpolation error bounds and
edge-concave underestimators, tree dynamic programming and min-marginals,
optimality-based bound tightening, rational height bounds for quadratic
programs, multiparametric quadratic programming, graph cuts and total
unimodularity. Section~\ref{sec:related} discusses this prior work in detail.
To the best of our knowledge, the following are new:
\begin{enumerate}[label=(\roman*)]
\item the node-separable form of the vertex lower bound, which makes the best
cellwise bound over a product grid computable by dynamic programming on a tree
decomposition, together with min-marginal filtering and the path certificate
(Section~\ref{sec:grids});
\item the accuracy-independent grid bound under global growth, with the
scale-invariant condition number $\bar\kappa$, and the resulting algorithms that
are fixed-parameter tractable in bag size and condition number for certified
approximation of sparse mixed-integer polynomial problems and exact solution
of sparse mixed-integer box quadratic programs, without knowledge of the growth
constant (Sections~\ref{sec:growth} and~\ref{sec:exact});
\item the integration of value functions, certified convex responses,
minimum-cut residual oracles and cut-based grid oracles into the
corrected-grid framework, with the corresponding curvature and growth transfer
(Section~\ref{sec:recourse});
\item the exactly feasible TU extension and the analysis of several
minimizers (Sections~\ref{sec:constraints} and~\ref{sec:optsets});
\item the lower bounds and limiting examples of Section~\ref{sec:limits},
in particular the expanding-box chain, the unique-minimizer hardness at width
three and the finite-order moment obstruction;
\item an exact-arithmetic implementation whose certificates are replayed by an
independent checker, and experiments confirming the predicted grid sizes
(Section~\ref{sec:computation}).
\end{enumerate}

\paragraph{Organization.}
Section~\ref{sec:related} reviews related work and Section~\ref{sec:setting}
fixes the model and the parameters. Sections~\ref{sec:grids}--\ref{sec:exact}
develop the corrected grids, the growth analysis and exact output.
Section~\ref{sec:recourse} treats conditional recourse,
Section~\ref{sec:constraints} coupling constraints and Section~\ref{sec:optsets}
several minimizers. Section~\ref{sec:limits} collects structural limits,
Section~\ref{sec:computation} reports the implementation and experiments, and
Section~\ref{sec:conclusion} lists open problems. Longer proofs are in the
appendices.
